\documentclass[11pt,a4paper]{article}
\usepackage[margin=25mm]{geometry}
\usepackage[T1]{fontenc}
\usepackage{lmodern,amsmath,amssymb,booktabs,longtable}
\usepackage{xurl}
\usepackage[hidelinks]{hyperref}
\usepackage{microtype}
\setlength{\emergencystretch}{2em}
\title{A reader's guide to the certificates\\for rational Fermat--Cantor exponents}
\author{Computational companion to \emph{Solutions of $x^\alpha+y^\alpha=1$ in $\lambda$-Cantor sets}}
\date{24 October 2026}
\begin{document}
\maketitle
\begin{abstract}
This guide explains the exact certificates used for the rational-exponent
part of Theorem 3.1. It collects the detailed reduction, parameter covers,
exceptional sources and verification predicates formerly printed in the
manuscript, and connects them to the executable checks. The certificate
data and algorithms are unchanged. The irrational-exponent proof and the
other current manuscript results are analytic.
\end{abstract}
\tableofcontents

\section{What the computation proves}
Let $C$ be the middle-third Cantor set and write
$E_{\alpha,1/3}:=\{(x,y)\in C^2:x^\alpha+y^\alpha=1\}$.
The required assertion is
\[
\dim_{\mathrm H}(E_{\alpha,1/3})>0
\qquad(\alpha\in\mathbb Q,\ \alpha>1).
\]
The case $\alpha=2$ has an independent analytic proof in manuscript
Theorem 2.11. For the remaining rational exponents, the computation either
supplies a strict logarithmic source in an admissible slope band or
verifies two strict returns on a specified Cantor rectangle. Both routes
produce an infinite construction; counting points at a fixed depth is
not the proof.

Read Section 2 for the reduction and coverage partition, Section 3 for the
exact inequalities checked by the programs, and Section 4 for replay
commands. A short first inspection is the sixteen-row denominator cover
in Section 3. The seven exceptional rational sources then illustrate why
the direct-return route is also needed.

\paragraph{Version and scope.}
This guide transcribes the rational proof and Appendix A from Overleaf
commit \nolinkurl{6e55513bbf5e706c570fdbbc634254b786a3b1d9}.
All numerical inputs and verification programs are the existing
\href{https://github.com/Nero-17/LEAN-Formalisation-solutions-of-fermat-cantor-equation/tree/v1.2.0}{v1.2.0 certificate release}.
The new guide does not report a new full replay of these calculations.
Previously completed checks and their logs are identified in the repository.
The manuscript's current Lemma 2.9, Theorem 2.11 and Theorem 4.2 use
analytic proofs; their older computational material remains available as
historical or supplementary material, not as a dependency of those proofs.
These are exact Python checks, not Lean formalizations.

\paragraph{Notation and analytic inputs.}
Program fields named \texttt{n} denote the manuscript exponent $\alpha$.
The admissible slope bands are the interval $(11/9,2027/1624)$ and
\[
(89/73,171/137),\quad(101/79,183/143),\quad(89/57,171/89),
\]
\[
(143/61,237/101),\quad(137/57,219/89),
\]
together with their reciprocals. Their eventual-periodicity property is
proved in manuscript Lemmas 2.5 and 2.9. The interval-cover principle,
strict-return principle and irrational-tangent upgrade are manuscript
Lemmas 2.3, 2.8 and 2.10, respectively. References to these numbered results
below are references to the manuscript.
For the normalization $(rs+t)/2$, the switch set is
\[
B_r:=\left[\frac13,\frac{r+1}{6}\right]\cup
\left[\frac r3,\frac{r+3}{6}\right]\cup
\left[\frac{r+1}{3},\frac{3r+1}{6}\right].
\]
Each accepted return record specifies the slope, initial level, return
depth and strict margin. Two words are incomparable when neither is a
prefix of the other. The strict-return lemma turns such a record and its
nonlinear error bound into the lower bound $\log2/(M\log3)$.

\section{Detailed reduction and coverage}

\emph{Strict logarithmic sources.}
Write \(\alpha=a/b>1\) in lowest terms. The case \(\alpha=2\) follows from
Theorem~2.11. For every other such exponent,
there are no solutions with both coordinates positive and rational.
To see this, suppose otherwise. In a normal closure, every conjugate
of \(x^{a/b}\) is \(x^{a/b}\) times a \(b\)-th root of unity, and
likewise for \(y^{a/b}\). The conjugate equation and equality in the
triangle inequality force both conjugates to remain positive real.
Thus \(x^{a/b},y^{a/b}\) are rational. Prime valuations and
\(\gcd(a,b)=1\) imply \(x=u^b,y=v^b\) for positive rationals
\(u,v\). Since \(a\geq3\), the equation \(u^a+v^a=1\) contradicts
Fermat's Last Theorem~(see the references in the manuscript).

Let \(u=(P+s)/3^p\), \(v=(Q+t)/3^q\), where \(P,Q>0\) have
padded ternary digits in \(\{0,2\}\). A strict logarithmic source
consists of \(d\in b^{-1}\mathbb Z\) satisfying
\begin{align*}
d\log3+\alpha\log(P/3^p)-\log \alpha-\log((Q+1)/3^q)&<0,\\
d\log3+\alpha\log((P+1)/3^p)-\log \alpha-\log(Q/3^q)&>0,
\end{align*}
with
\[
\left[\frac{\alpha Q}{P+1},\frac{\alpha(Q+1)}P\right]
\]
strictly inside an admissible slope interval or its reciprocal.
If \(d=j/b\), choose arbitrarily large integers \(k\) satisfying
\(ak+j\equiv0\pmod b\). Then \(\alpha k+d\) is an integer and
\[
x=3^{-k}u,\qquad y=1-3^{-\alpha k-d}v
\]
are Cantor cylinders. The exact logarithmic residual
\[
d\log3+\alpha\log u-
\log\!\left(3^{\alpha k+d}[1-(1-3^{-\alpha k-d}v)^\alpha]\right)
\]
converges in \(C^1\) to \(d\log3+\alpha\log u-\log \alpha-\log v\).
For real powers the required identity is
\[
\frac{1-(1-z)^\alpha}{z}=\alpha\int_0^1(1-tz)^{\alpha-1}\,dt.
\]
Thus the strict corner and slope inequalities persist. Reflecting
the \(v\) coordinate makes both partial derivatives positive;
Lemma~2.3 gives a nonempty compact fibre.
If it had an isolated point, the vanishing-error digital argument
of Lemmas~2.5 and~2.9, as in
Theorem~2.6, would make both coordinates rational.
The preceding obstruction excludes this. The fibre is perfect,
and its tangent ratio is strictly monotone along the Fermat arc.
Lemma~2.10 gives positive dimension.

\emph{Coverage of the rational exponents.}
The interval of permitted \(d\) has length
\[
\frac{\alpha\log(1+1/P)+\log(1+1/Q)}{\log3}.
\]
For denominators \(b\geq110\), a uniform construction makes this
length greater than \(1/110\). Write \(\alpha=3^m t\), \(1\leq t<3\).
The sixteen rows \((P,Q,\ell,u)\) in
Section~\ref{app:real-certificates} satisfy the following exact
open-cover condition:
\[
[1,3]\subset\bigcup
\left(\max\left\{\frac{\ell(P+1)}Q,
\left(\frac{11}{1100}-\frac1{Q+1}\right)(P+1)\right\},
\frac{uP}{Q+1}\right).
\]
Replace \(P\) by \(3^mP\), appending zero ternary digits.
The derivative interval stays strictly in \((\ell,u)\), and
\[
\alpha\log\left(1+\frac1{3^mP}\right)
+\log\left(1+\frac1Q\right)
>\frac{t}{P+1}+\frac1{Q+1}>\frac{11}{1100}.
\]
Since \(\log3<11/10\), the permitted interval contains a point of
\(b^{-1}\mathbb Z\) whenever \(b\geq110\).

Integer exponents are covered by the certificates and nine direct
returns detailed in Section~\ref{app:integer-certificates}.
For \(2\leq b\leq109\), exact verification gives the following
exhaustive decomposition:
\[
\begin{array}{c|l}
\text{exponent range}&\text{coverage}\\ \hline
1<a/b<2187&7{,}738{,}063\text{ width sources},\\
 &295{,}480\text{ individual sources, seven direct returns}\\
3^7\leq \alpha\leq3^8&650\text{ closed parameter cells}\\
3^8\leq \alpha\leq3^9&508\text{ cells}\\
3^9\leq \alpha\leq3^{10}&470\text{ cells}\\
3^{10}\leq \alpha\leq3^{11}&552\text{ cells}\\
3^{11}\leq \alpha\leq3^{12}&547\text{ cells}\\
\alpha\geq3^{12}&\text{continuous covers and uniform tail}.
\end{array}
\]
There are exactly \(8{,}033{,}550\) reduced fractions in the first
range. The five closed covers verify every denominator separately
and have no gaps at their joins. The continuous covers for
\(\alpha\geq3^{12}\), described in Section~\ref{app:integer-certificates},
cover every real normalized exponent \(t\), with integer \(d\).
For rational \(\alpha=a/b\), taking \(k\) divisible by \(b\) makes the
same transfer valid. Thus their use does not require integrality
of \(\alpha\).

The seven remaining exponents have direct-return certificates:
\[
\begin{array}{c|rrrrrrr}
\alpha&24/23&19/18&13/12&8/7&8/3&11/4&53/2\\ \hline
M&19&17&17&15&17&17&16.
\end{array}
\]
Each certificate specifies a genuine Cantor rectangle, a rational
slope \(r\), an initial level \(T\in B_r\), and a normalization
\((rs+t)/2-T+H(s,t)\), with \(H(0,0)=0\).
It verifies two incomparable returns from every point of \(B_r\),
of length at most \(M\), with margin \(\eta\), and
\[
\sup_{[0,1]^2}(|\partial_sH|+|\partial_tH|)
<\frac{\eta}{2\,3^M}.
\]
Here \(\eta=10^{-7}\) for \(8/3\) and \(\eta=10^{-5}\) otherwise.
The exact recursion and algebraic error bounds are given in
Section~\ref{app:real-certificates}. Lemma~2.8
then gives the lower bound \(\log2/(M\log3)>0\) in each case.
This proves the rational case and completes the proof of
Theorem~3.1.


\section{Exact verification and certificates}\label{app:certificates}

The computations below use exact integer or rational inequalities.
We retain their verification predicates and mathematical implications
here. The certificates, verification scripts, calculation
records and reproduction instructions are available in the
\href{https://github.com/Nero-17/LEAN-Formalisation-solutions-of-fermat-cantor-equation/tree/v1.2.0}
{versioned computational companion repository}.
These are Python certificate checks; they are not Lean formalizations.

\subsection{Ordinary binary returns}
The underlying exact recursion can be described without a grid.
Let \(K\) be the union of the three closed switch intervals
shrunk by the specified margin. Starting with \(A_0=K\) and
\(D_0=\varnothing\), set
\begin{align*}
A_{j+1}&=A_j\cup\bigcup_c(A_j+c)/3,\\
D_{j+1}&=D_j\cup\bigcup_c(D_j+c)/3\\
&\quad{}\cup\bigcup_{c\ne d}
       \bigl((A_j+c)/3\cap(A_j+d)/3\bigr),
\end{align*}
where \(c,d\in\{0,1,r,r+1\}\).
The set \(A_j\) records one return of length at most \(j\);
\(D_j\) records two incomparable returns. Induction proves
both interpretations, since the last union chooses different
first digits. Exact unions and intersections of closed rational
intervals implement the recursion. The assertion
\(B_r\subset D_{L_\alpha}\), together with the nonlinear error bound
in the proof, is the required finite certificate.

\subsection{Logarithmic sources and the infinite exponent tail}
For each logarithmic source we verify the ternary digits, the two
strict logarithmic corner signs, and the full rational
derivative-ratio interval.
We enclose logarithms by exact range reduction and
\[
\log z=2\sum_{j=0}^{J-1}
  \frac{((z-1)/(z+1))^{2j+1}}{2j+1}+R_J,\qquad z\geq1,
\]
where
\[
0\leq R_J\leq
\frac{2((z-1)/(z+1))^{2J+1}}
 {(2J+1)(1-((z-1)/(z+1))^2)} .
\]
For \(z<1\), use \(\log z=-\log(1/z)\).
Range reduction uses exact powers of two and the same enclosed
value of \(\log2\). All endpoint rounding is outward.

The fixed-scale and tail checks test each closed parameter cell,
each Cantor rectangle, and its assigned slope band.
They compute inner corner-image intervals
with a strict inset \(10^{-12}\), verify that their translates
cover the circle, and then verify that the closed parameter
cells cover \([1,3]\).
For the tail, the integral enclosure printed in the proof
controls every \(Z\geq1180980\); testing a finite list of
exponents is not used in its place.

\subsection{Integer sources and the continuous exponent tail}\label{app:integer-certificates}
The logarithmic source criterion in Theorem~3.1 also applies when the reduced denominator is one. The circle case is supplied by Theorem~2.11.

Here is the finite covering argument that supplies these strict
sources. For large exponents write \(\alpha=3^m t\), \(1\leq t<3\).
Choose a Cantor endpoint \(u_0=P/3^D\), an integer \(h\), and
\[
u=u_0+3^{-m}a,\qquad a\in3^hC,\qquad m\geq D+h .
\]
The depth condition makes every such \(u\) a Cantor point.
Modulo one, the logarithmic source condition asks whether
\[
\frac{\log v-3^mt\log(1+a/(3^mu_0))}{\log3}
\]
contains the target phase
\((\alpha\log u_0-\log \alpha)/\log3\).
It suffices that strict inner corner images of admissible
rectangles cover the entire circle. Adding a constant rotation,
we check instead the phase
\[
\frac{\log v+t3^m
        \log((3^mu_0+3^h/2)/(3^mu_0+a))}{\log3}.
\]
At fixed \(m\), its corner values are affine in \(t\), so their
bounds on a closed parameter interval occur at its endpoints.
For a uniform tail in \(m\), use the exact identity
\[
Z\log\frac{Z+3^h/2}{Z+a}
       =\int_a^{3^h/2}\frac{Z}{Z+s}\,ds,\qquad Z=3^mu_0.
\]
For fixed \(a\), this is monotone in \(Z\), with direction
determined by the orientation of the integral. Its value for
every \(Z\geq Z_{\min}\) is enclosed by its value at \(Z_{\min}\)
and its limit \(3^h/2-a\). Derivative ratios have analogous
rational endpoint bounds. Thus finitely many closed coefficient
cells certify the whole infinite tail.

The first two rows below concern integer exponents; the remaining rows cover continuous exponent ranges. The exact certificate decomposition is
\[
\begin{array}{c|l}
\text{exponents}&\text{verification}\\ \hline
2\leq \alpha\leq300&\text{strict sources and the nine returns below}\\
301\leq \alpha\leq531440&\text{individual strict logarithmic sources}\\
3^{12}\leq \alpha<3^{13}&8921\text{ closed parameter cells}\\
3^{13}\leq \alpha<3^{14}&6027\text{ closed parameter cells}\\
3^{14}\leq \alpha<3^{15}&5310\text{ closed parameter cells}\\
\alpha\geq3^{15}&5862\text{ coefficient cells, }293\text{ endpoint choices}.
\end{array}
\]
For the three fixed scales \(D=9,h=3\), and the two parameter
domains \([1,2]\), \([2,3]\) are covered without gaps.
For the tail, \(D=11,h=3,u_0\geq20/243\), and
\(Z_{\min}=1180980\); the symbolic condition \(15\geq11+3\)
holds. The individual-source range ends at \(3^{12}-1\).
Thus these certificates leave no uncovered integer.
Section~\ref{app:certificates} specifies their data, exact
logarithm bounds, and verification predicates.

Nine small exponents use direct returns instead of admissible
slopes:
\[
\begin{array}{c|rrrrrrrrr}
\alpha&9&10&23&25&28&31&32&78&82\\ \hline
L_\alpha&21&15&16&17&18&16&20&16&17 .
\end{array}
\]
For completeness, the common transfer is as follows. In a
certified cylinder \(x=(P+b)/3^p,\ y=(Q+a)/3^q\), \(q\geq p\),
multiply the Fermat residual by \(3^{\alpha q}\) and divide by
\(2\alpha3^{\alpha(q-p)}P^{\alpha-1}\). It becomes
\[
\frac{ra+b}{2}-T+H(a,b)=0,\quad
r=\frac{Q^{\alpha-1}}{3^{\alpha(q-p)}P^{\alpha-1}},\quad
T=\frac{3^{\alpha q}-Q^\alpha-3^{\alpha(q-p)}P^\alpha}
         {2\alpha3^{\alpha(q-p)}P^{\alpha-1}}.
\]
The remainder vanishes at zero and has sup-norm Lipschitz constant
at most
\[
\frac{(Q+1)^{\alpha-1}-Q^{\alpha-1}}{2\,3^{\alpha(q-p)}P^{\alpha-1}}
+\frac{(P+1)^{\alpha-1}-P^{\alpha-1}}{2P^{\alpha-1}}.
\]
If the certificate uses a smaller rational slope \(r_0\), add
\((r-r_0)/2\) to this bound. The exact checks give \(T\in B_{r_0}\),
two incomparable returns of length at most \(L_\alpha\) into the
\(1/100000\)-interior of \(B_{r_0}\), and the strict inequality
\(3^{L_\alpha}\operatorname{Lip}(H)<1/100000\), including this slope
adjustment. Descendant normalization preserves the same error
bound. Lemma~2.8 therefore gives
\(\dim_{\mathrm H}(E_{\alpha,1/3})\geq\log2/(L_\alpha\log3)>0\).
Together with the strict sources, this proves the integer case used in Theorem~3.1.


\subsection{Rational-exponent sources and coverage}\label{app:real-certificates}
The data and exact checkers for Theorem~3.1
are included in release v1.2.0 of the repository cited above.
The guide \nolinkurl{docs/real-exponents.md} identifies the accepted
inputs, verification commands and previously completed verification
records; \nolinkurl{docs/paper-index.md} connects the paper's results
to their computational files. All acceptance tests use integer or rational
arithmetic, including outward enclosures of logarithms and algebraic
roots; finite-depth growth estimates play no role in the proof.

The sixteen-source denominator cover uses the following rows.
Every \(P,Q\) has ternary digits in \(\{0,2\}\). Substituting each
row in the open-interval formula in the proof and sorting the resulting
rational endpoints verifies strict overlaps covering \([1,3]\).
All listed slope bands are admissible bands or their reciprocals.
\[
\begin{array}{rr|cc}
P&Q&\ell&u\\ \hline
180&222&89/73&171/137\\
224&180&137/171&73/89\\
180&218&89/73&171/137\\
182&218&89/73&171/137\\
216&168&137/171&73/89\\
186&218&89/73&171/137\\
162&242&89/57&171/89\\
182&224&89/57&171/89\\
186&188&89/57&171/89\\
224&186&89/57&171/89\\
240&164&89/57&171/89\\
242&162&89/57&171/89\\
504&216&89/73&171/137\\
512&216&89/73&171/137\\
222&182&137/57&219/89\\
540&222&89/73&171/137\\
\end{array}
\]


For the finite range \(1<a/b<2187\), \(2\leq b\leq109\),
the checker enumerates every coprime pair. It first tests the
denominator-dependent width cover, then checks an individual source
where necessary. Each individual record gives \((a,b,P,Q,J)\) and
the slope band. With \(p,q\) the ternary depths, set
\(d=(J+ap-bq)/b\). The corner signs are exactly
\[
3^Jb^bP^a<a^b(Q+1)^b,\qquad
3^Jb^b(P+1)^a>a^bQ^b.
\]
Negative powers of three are cleared before an integer comparison.
For the large collection, range-reduced logarithms are enclosed by
the positive series for \(2\operatorname{arctanh}z\), \(0\leq z\leq1/3\),
with explicit tail bounds and outward integer rounding. A separate
implementation uses seventy terms and 200-bit fixed-point arithmetic.
It re-enumerates the entire domain and verifies the partition
\[
8{,}033{,}550=7{,}738{,}063+294{,}016+1{,}464+7.
\]
The \(1{,}464\) additional sources are also checked by direct integer
powers. The seven remaining pairs are precisely those in the
direct-return table in the proof.

For the five layers \(m=7,\ldots,11\), write \(\alpha=3^mt\),
\(1\leq t\leq3\), and use the rotated phase displayed in
Section~\ref{app:integer-certificates}. At each fixed endpoint
\(u_0=P/3^D\), the admissible rectangles have strict derivative
bounds throughout a closed \(t\)-cell. Their corner-image intervals
are shrunk by \(10^{-12}\). For each assigned denominator \(b\),
multiply these intervals by \(b\) and reduce modulo one. Exact
interval unions cover the whole circle, which proves coverage of
every target phase modulo \(1/b\). Different denominators may use
different endpoints, but every assigned endpoint by itself covers
its denominator's entire phase circle. Intervals from distinct
endpoints are never combined without their rotations. The checker
reconstructs the corner and derivative bounds and verifies all
closed-cell joins from \(1\) to \(3\).

Six special exponents use the following positive source cylinders
and actual integer scalings:
\[
x=3^{-k}\frac{P+s}{3^p},\qquad
y=1-3^{-m}\frac{Q+1-t}{3^q},\qquad s,t\in C.
\]
\begin{center}
\small
\begin{tabular}{c|r|r|c}
\(\alpha\)&\(P\)&\(Q\)&\((p,q,k,m)\)\\ \hline
\(24/23\)&9006915410325491654&10818090212415614952&\((40,40,110,115)\)\\
\(19/18\)&9110013456467818106&10841815821954030882&\((40,40,104,110)\)\\
\(13/12\)&9120707925586079678&10817949568504999626&\((40,40,105,114)\)\\
\(8/7\)&9450592763918769000&10918429022892060476&\((40,40,68,78)\)\\
\(11/4\)&48241792345204752&37518532384253874&\((35,35,103,284)\)\\
\(53/2\)&105672039117920552328&8521470282198373850&\((42,40,101,2680)\)\\
\end{tabular}
\end{center}

In each row \(d=m-\alpha k\in b^{-1}\mathbb Z\).
For the limiting equation \(u^\alpha-\alpha3^{-d}v=0\), divide by
\(2\alpha3^{-d-q}\). Its tangent ratio \(r_*\) at the lower corner
and its level satisfy
\[
r_*^b=3^{bd+b(q-p)}(P/3^p)^{a-b},\qquad
T=\frac{Q+1-(b/a)r_*P}{2}.
\]
The certificate encloses these quantities by rational root brackets,
selects a rational reference slope, and bounds the full variation
of both derivatives. For \(1<\alpha<2\),
\((1-z)^{\alpha-2}\leq3/2\) on the actual cylinders; for \(\alpha\geq2\)
it is at most one. Taylor's theorem therefore bounds the additional
normalized constant and \(t\)-derivative errors, respectively, by
\[
\frac{3(\alpha-1)}8\,3^{q-m},\quad
\frac{3(\alpha-1)}4\,3^{-m}
\quad(1<\alpha<2),
\]
and by the same expressions with \(3/2\) replaced by \(1\)
when \(\alpha\geq2\). The checks include the shifted actual level,
the reference-slope error and the strict derivative margin.

The seventh source, \(\alpha=8/3\), is directly on the actual curve
and is specified in the accompanying certificate. Rational brackets for the
cube roots of the fifth powers of its two coordinate endpoints
enclose both first derivatives; multiplication by the endpoints
also encloses the residual. The two source checks use independent
root-enclosure implementations.

For all seven sources, the binary-return recursion given above
is applied to the \(\eta\)-interior of \(B_r\), using the specified
rational reference slope. It verifies the full inclusion
\(B_r\subset D_M\), not a sample of levels. Together with the
strict nonlinear error bound in the theorem, this proves that
the two incomparable returns can be repeated at every descendant.
For \(8/3\), a second interval implementation reproduces all
seventeen generations of the recursion, ending with \(289{,}323\)
one-return intervals and \(1{,}787{,}353\) two-return intervals.



\section{How to replay the verification}
The \href{https://github.com/Nero-17/LEAN-Formalisation-solutions-of-fermat-cantor-equation/blob/v1.2.0/docs/real-exponents.md}{all-real certificate guide}
lists the individual input files and direct commands. The
\href{https://github.com/Nero-17/LEAN-Formalisation-solutions-of-fermat-cantor-equation/blob/v1.2.0/docs/verification.md}{verification guide}
describes the retained integer stages. Use Python 3.11 or later, with its
standard library, without \texttt{-O}. Run from the repository root:
\begin{verbatim}
git checkout v1.2.0
python reproduce.py audit
python reproduce.py returns
python reproduce.py integer-sources
python reproduce.py phase-covers
python reproduce.py assembly
python reproduce.py real-quick
python reproduce.py real-sources
python reproduce.py real-phase-covers
python reproduce.py real-returns
python reproduce.py real-assembly
\end{verbatim}
The \texttt{returns} suite also replays an older weighted circle bound;
that extra check is not needed for the current analytic Theorem 2.11.
The \texttt{uniform-bound} suite concerns the older proof of Theorem 4.2
and is not needed for the current manuscript. The broader \texttt{all}
suite additionally includes these historical checks and numerical-figure
material. A full replay may take substantial time and memory.

\begin{center}
\begin{tabular}{p{.27\linewidth}p{.65\linewidth}}
\toprule
Stage & Mathematical task \\
\midrule
\texttt{audit} & Check input identity against stored hashes; not proof verification. \\
\texttt{returns} & Small-integer sources, nine direct returns and source transfer. \\
\texttt{integer-sources} & Individual integer sources from 301 to 531440. \\
\texttt{phase-covers} & Continuous exponent scales and the infinite tail. \\
\texttt{real-quick} & Sixteen-row denominator cover and two source checks at $8/3$. \\
\texttt{real-sources} & Low-denominator source checks and exceptional sources. \\
\texttt{real-phase-covers} & The five intermediate scales for denominators 2 to 109. \\
\texttt{real-returns} & The exceptional rational return recursions. \\
\texttt{assembly}, \texttt{real-assembly} & Coverage joins and receipt consistency after the substantive checks. \\
\bottomrule
\end{tabular}
\end{center}

The driver expands compressed inputs in a separate workspace. Use
\texttt{python reproduce.py all --list} to inspect the underlying commands.
An assembly check or a saved success flag does not replace the arithmetic
checks. Likewise, the large-exponent adoption check
\nolinkurl{thirdreal20261021_adversarial_large.py} checks hashes and joins,
not the original logarithm calculations: the \texttt{phase-covers} stage
performs those calculations.

\paragraph{What to inspect in a certificate.}
Check that the endpoints have legal Cantor digits; that strict corner
signs and derivative intervals hold on the whole assigned domain; that
all parameter cells join without gaps; and, for direct returns, that the
whole switch set is included in the two-return set. The nonlinear error
must be smaller than the return margin after normalization. These are
the finite predicates used by the analytic transfer. Search programs and
failed candidate files are discovery history, not accepted certificates.

\paragraph{Recorded evidence.}
The repository retains \nolinkurl{docs/source-manifest.json},
\nolinkurl{docs/reproduction-results.json},
\nolinkurl{docs/reproduction-results-v1.1.0.json},
\nolinkurl{docs/reproduction-results-v1.2.0.json}, and the corresponding
\texttt{records/} directories. They identify prior checks and provenance.
This document relocates and explains that evidence; it does not add a new
mathematical claim or report computations that were not rerun.
\end{document}
